\section{Conclusions}
\label{sec:conclusion}

This study examined the dependency structure of Brazilian equities traded on B3 through an integrated econophysics framework combining correlation analysis, Random Matrix Theory, hierarchical representations, filtered financial networks, and a retrospective volatility-model assessment. The empirical evidence shows that the Brazilian equity market exhibits positive but heterogeneous cross-sectional dependence, with stronger relationships within economic sectors and substantial variation in market-wide synchronization over time.

The spectral analysis shows that this dependence is strongly influenced by a dominant collective market component, while additional significant modes retain more localized patterns that do not map directly onto simple sector classifications. Removing the leading market mode changes the geometry of the dependency structure rather than eliminating it. Hierarchical clustering, MST, and PMFG representations show that weaker but organized relationships remain, with changes in local clustering, network connectivity, and the identity of central assets. The subsector-level representation further translates these residual dependencies into an economically interpretable view of the market.

In the validation-selected Ridge comparisons, Set B and Set C have higher test QLIKE than Set A at both horizons, and the paired HAC tests favor the Set A baseline after adjustment across the four structural contrasts. HAR-RV also provides an explicit benchmark; its paired comparison with Ridge Set A is not statistically distinguishable at the 5% level. The neural sensitivity analysis shows materially greater seed variation for MLP than for CNN. Because the structural descriptors use the full synchronized sample, these findings are retrospective and do not establish a point-in-time forecasting effect.

The study integrates spectral filtering and network representations to examine broad market synchronization and localized dependency structure in Brazilian equities. A next step is to reconstruct RMT and network descriptors recursively at each forecast origin and test them using only information available at that time.
